\chapter{Реализация}
\label{chap:impl}

\section{Архитектура языкового сервера Jasmin}

В этом разделе представлена реализация языкового сервера (LSP) для Jasmin,
разработанного на OCaml. Сервер напрямую интегрируется с фронтендом компилятора
Jasmin, используя его систему типов и абстрактное синтаксическое дерево (AST)
для предоставления функциональности через протокол LSP.

Ключевые части логики реализованы в виде чистых функций, что упрощает отладку и
расширение. Сервер работает как автономный исполняемый файл и управляется через
конечный автомат состояний ($Uninitialized$, $Running$, $Stopped$), асинхронно
обрабатывая запросы клиента в отдельных потоках \texttt{lwt}. Реализовано
конфигурируемое логирование для удобства отладки.

\section{Реализованные функции}

\label{impl:code-diagnostics}

\subsubsection*{Диагностика кода в реальном времени} 

Диагностика запускается при сохранении файла и позволяет выявлять синтаксические
и типовые ошибки без вызова компилятора. Процесс выполняется в отдельном потоке
и состоит из двух этапов: сначала парсинг кода, затем — проверка типов с помощью
API компилятора. Текущая реализация не поддерживает инкрементальный анализ и
обрабатывает файл целиком.

\subsubsection*{Подсветка символов и навигация}

Для подсветки всех использований символа сервер находит соответствующий узел в
AST и ищет все его вхождения, используя уникальные идентификаторы, которые
компилятор присваивает каждому символу. Это позволяет корректно различать
переменные с одинаковыми именами.

\subsubsection*{Переход к определению}
Функция перехода к определению (`Go to definition`) также реализована с помощью
AST, поскольку узлы дерева уже содержат информацию о местоположении определения
символа.

\subsubsection*{Отображение типа при наведении} 

Для отображения типа символа при наведении курсора используется API компилятора,
а именно его функции для форматированного вывода (pretty-printing). Это
позволяет получать и отображать информацию о типе напрямую из AST без
дополнительных преобразований.

\section{Поддержка Tree-sitter}

\label{impl:tree-sitter}

Для подсветки синтаксиса был разработан парсер на основе Tree-sitter
\cite{tree-sitter}. Грамматика была адаптирована из существующей грамматики для
Menhir~\cite{menhir}, используемой в компиляторе Jasmin.

Tree-sitter строит конкретное синтаксическое дерево (CST) кода и использует
специальные запросы (`highlight queries`) для раскраски синтаксических
конструкций. В отличие от семантической подсветки, этот механизм основан только
на структуре кода, но значительно улучшает читаемость и удобство работы в
редакторе.